% Section 6, Related work (about 700 words; up to ~1,200 while drafting): problem-space attacks, constraints over
% features, robust and monotone detectors, incremental computation and provenance. Abstract interpretation is cited
% here and in Sound Directions (R6, 2026-11-12; notes in docs/evidence/r6-prior-work.md).
% Spine: docs/writing-spec.md reaction log (2026-11-03).

\paragraph{Problem-space attacks}
Our work builds on the formalization of problem-space attacks by Pierazzi et al.~\cite{pierazzi2020intriguing}. They
describe an attack as a sequence of transformations applied to an object under problem-space constraints, and relate the
result to the feature space through the feature mapping, side-effect features included. For feature mappings that are
neither differentiable nor invertible, those side effects generally cannot be computed in
advance~\cite{pierazzi2020intriguing}. Šrndić and Laskov, evading PDFrate~\cite{srndic2014practical}, found it
infeasible to enumerate the dependencies among its features and account for their effects a priori, and accepted files
whose features differed from those their attack had planned. The toolkits now used against Windows malware
detectors~\cite{demetrio2021adversarial,demetrio2021functionality,song2022mab} instead find the side effects by applying
a transformation and extracting the features of the result. Such an
attack learns only about the edits its search reaches, and if it fails, other edits may still evade the detector. For
sum-then-readout extractors (Definition~\ref{def:fw:extractor}), our work computes the side effects from the extractor
itself: exactly for a given edit and, for every edit in a declared set, as the direction in which each feature can move,
including edits no attack has tried.

% Paragraphs 1-4 revised after the user's written review (2026-11-04) and accepted as a draft; the user will fine-tune
% them by hand.
\paragraph{Constraints over features}
A second line of work keeps the attack in feature space, where gradients and search are cheap, and restricts it to
realizable feature vectors by constraints over features: some features are immutable, and others are tied together, as
a total must equal the sum of its parts. Simonetto et al.~\cite{simonetto2022unified} made this approach
domain-independent, with a grammar for linear and non-linear relations among features and attacks that respect any
constraint set written in it, and Chernikova and Oprea's FENCE~\cite{chernikova2022fence} does the same for
network-traffic and malicious-domain classifiers, updating dependent features after each step. Kantchelian et al.~\cite{kantchelian2016evasion} find the nearest evading instance of a tree ensemble with a
mixed integer program, to which the analyst can add dependencies that the extractor creates, such as features of which
exactly one is nonzero.

All three take the constraints as input, which leaves open where the constraints come from. Dyrmishi et
al.~\cite{dyrmishi2023empirical} note that they must be written in analytical form, typically with a domain expert; for
Windows malware, they took the immutable features from Microsoft's description of the executable format and the
dependencies among features from sample files. Sheatsley et al.~\cite{sheatsley2021robustness} learn constraints from
data as Boolean formulas, and show that a learned constraint set admitting only valid inputs reduces to memorizing the
training data. Two works learn constraints to generate examples for adversarial training: Bostani et
al.~\cite{bostani2025level} learn dependencies among features from many Android apps, and note that constraints learned
from a finite sample describe only part of the realizable feature space; Labaca-Castro et
al.~\cite{labaca2022realizable} estimate, for each EMBER feature, how likely a problem-space
attack is to change it, and sample perturbations from these estimates as if the features were independent. None of
these derives the constraints from the extractor itself. We do: for a declared set of edits, we compute which way each
feature can move under any edit in the set, with neither an expert nor a sample.

\paragraph{Hardening by retraining}
The knowledge that constrains an attack can also harden a detector, by retraining it on adversarial examples, but only
if the examples are realizable: Dyrmishi et al.~\cite{dyrmishi2023empirical} found that hardening a malware detector on
unconstrained feature-space examples brought little robustness against a problem-space attack, because those examples
moved the features in other directions than the problem-space ones did. Realizable examples are costly to generate.
Lucas et al.\ made adversarial training of raw-binary detectors practical by generating them in parallel on over 140
workers~\cite{lucas2023adversarial}, and later trained robust detectors without generating any adversarial
examples~\cite{lucas2024training}. Tsingenopoulos et al.~\cite{tsingenopoulos2024train} generate examples with a
reinforcement-learning agent that applies problem-space transformations to dynamic-analysis reports, and measure
robustness as the probability that the strongest adversary with these transformations fails, estimated by training a
new agent on held-out reports. The transformations, and how a change to one entry of a report carries over to the
others, are written by hand; the second is what our framework computes.

\paragraph{Monotone and verified detectors}
Rather than resist the edits it was trained on, a detector can be built so that no edit evades it. If the edits can
only raise each feature they touch, and the detector's score never falls when a feature rises, no sequence of edits
lowers the score. Incer et al.~\cite{incer2018adversarially} built malware detectors on this principle: they decided by
hand, for each feature, whether an attacker can cheaply raise it, lower it, both, or neither, dropped the features that
move both ways, and trained a classifier monotone in the rest. Gimenez et al.~\cite{gimenez2025certifiably} proved the
converse, that every robust detector can be written as a monotone classifier over some post-processing of its features; they pick the
40 features of their certified detector by hand, as hard for an attacker to change, and learn the post-processing of
their empirically robust ERDALT from observed attacks. Chen et al.\ specify what an attacker can afford from domain
knowledge, training tree ensembles against costs that differ across features and between raising and lowering
one~\cite{chen2021cost}, and classifiers that satisfy monotonicity and other properties~\cite{chen2021learning}.
% Verified 2026-11-04 against arXiv 2105.11363: Property 1 and Sec. 6.1 ("Monotonic Directions": suspiciousness and
% economic cost); datasets: cryptojacking, Twitter spam accounts, Twitter spam URLs.
Their verifiably robust PDF classifiers~\cite{chen2020training} rely instead on the feature set: Hidost's
features~\cite{srndic2016hidost} are simple enough to show which of them inserting or deleting a subtree of the document
can change. Wang et al.~\cite{wang2022robust} also guarantee robustness by design, normalizing each input to a canonical
form under a relation they specify, for instance by merging groups of equivalent API calls listed with expert
knowledge. Each of these detectors thus needs to know what an attacker's edits can do to each feature, and obtains it
by hand, from observed attacks, or by choosing features that are easy to analyze, which excludes the feature sets
deployed detectors use, such as EMBER's. We compute it for every feature of a sum-then-readout extractor and any
declared set of edits, so each construction can be applied to such feature sets.

\paragraph{Certified defenses}
Certified defenses bound the attacker by how far its edits can move the input, whatever the edits mean.
RS-Del~\cite{huang2023rs} certifies a byte-sequence classifier such as MalConv against an adversary who inserts,
deletes, or substitutes up to a given number of bytes, by classifying many randomly thinned copies of the input, and
DRSM~\cite{saha2024drsm} certifies a MalConv variant against an adversary whose bytes are contiguous, by classifying
windows of the input separately. Giri et al.~\cite{giri2026towards} smooth a LightGBM classifier on EMBER's features,
ablating groups of features and adding noise, and certify a radius in feature space. These guarantees bound how much the
attacker may edit, up to a radius in bytes (RS-Del reaches a certified accuracy of 91\% at 128 bytes) or in feature
distance, and RS-Del's and Giri et al.'s hold only with high probability. Ours bounds what an edit can do, so it holds
for every edit of the threat model, of any size, with no sampling; the price is that it covers only the kinds of edit
the threat model declares, where theirs cover any edit within the radius.
% TODO cite (R21): Gibert et al.'s certified detectors (AISec 2023, IEEE Access 2024, as listed by Giri et al.) are not in
% papers/; papers/gibert2024assessing.pdf is the packing study (main.bib: gibert2025assessing), not these.

\paragraph{Incremental computation and provenance}
Our side-effect computation uses machinery from databases and programming languages, fields that have long asked how a
change to an input affects an output without recomputing the output. Incremental view maintenance updates an aggregate
over a table after rows are inserted or deleted~\cite{gupta1995maintenance,koch2014dbtoaster}. Its update rules for sums
and maxima are the rules by which an edit changes our sums and pools. Provenance semirings~\cite{green2007provenance}
annotate each output of a query with the inputs that produced it, which is what our item table records for each
feature. Self-adjusting computation~\cite{acar2002adaptive} and incremental parsing~\cite{wagner1998efficient} re-run
only the part of a program or a parse that a change can reach, and incremental parsing reparses only as far around an
edit as the parser's lookahead reaches, much as a write to a file can change only the scanned items around it, such as
the strings and windows it overlaps, unless it lands on a position the parser reads to find other items. Each of these
techniques updates the output after one edit that has been made. A defender does not know the attacker's edit, and asks
instead about every edit the threat model allows, on every input of a population, before any content is chosen.
Reasoning soundly about all executions of a program at once is the subject of abstract
interpretation~\cite{cousot1977abstract}. Calzavara et al.~\cite{calzavara2020certifying} use it to certify decision
trees against an attacker specified as rewriting rules that change individual features. Our framework applies this kind
of reasoning to the extractor rather than to the model. It bounds which way each feature can move under every edit of a
declared set, including writes to positions that locate items far from the write, and the bounds it returns are the
kind of feature-space description that such certification takes as input.
